\documentclass[12pt,a4paper]{article}
\usepackage[margin=25mm, headheight=15pt, headsep=10pt]{geometry}
\usepackage{fontspec}
\usepackage[table]{xcolor} % Note: table option is required for rowcolors
\usepackage{titlesec}
\usepackage{tocloft}
\usepackage{fancyhdr}
\usepackage{booktabs}
\usepackage{tcolorbox}
\tcbuselibrary{skins, breakable}
\usepackage[colorlinks=true, linkcolor=emerald, urlcolor=emerald, bookmarks=true]{hyperref}

\definecolor{emerald}{RGB}{0,102,51}
\definecolor{darkred}{RGB}{139,0,0}
\definecolor{lightgray}{RGB}{245,245,245}

\usepackage[nil]{babel}
\babelprovide[import, main]{bengali}
\babelprovide[import]{arabic}
\babelfont{rm}[Renderer=HarfBuzz,Scale=1.0]{Noto Serif Bengali}
\babelfont{sf}[Renderer=HarfBuzz]{Noto Sans Bengali}
\babelfont{tt}[Renderer=HarfBuzz]{Noto Sans Bengali}
\babelfont[arabic]{rm}[Renderer=HarfBuzz,Scale=1.1]{Noto Naskh Arabic}

% Section styling
\titleformat{\section}{\Large\bfseries\color{emerald}}{\thesection.}{0.5em}{}
\titleformat{\subsection}{\large\bfseries\color{emerald!85!black}}{\thesubsection.}{0.5em}{}
\titleformat{\subsubsection}{\normalsize\bfseries\color{emerald!70!black}}{\thesubsubsection.}{0.5em}{}
\titlespacing*{\section}{0pt}{18pt}{8pt}

% Fancy Header/Footer setup
\pagestyle{fancy}
\fancyhf{} % clear all header and footer fields
\fancyhead[L]{\color{emerald}\small\textbf{\nouppercase{\leftmark}}}
\fancyfoot[L]{\scriptsize\color{black!50}{\lat{AI}}-সংকলিত, আলেম-যাচাইকৃত নয়}
\fancyfoot[C]{\color{emerald!70!black}\thepage}
\renewcommand{\headrulewidth}{0.4pt}
\renewcommand{\headrule}{\hbox to\headwidth{\color{emerald}\leaders\hrule height \headrulewidth\hfill}}
\renewcommand{\footrulewidth}{0pt}

% Custom boxes
% 1. Ayat/Hadith Box
\newtcolorbox{ayatbox}[1][]{
    enhanced, breakable, 
    colback=emerald!5, 
    colframe=emerald, 
    boxrule=0.5pt, 
    arc=3pt, 
    left=10pt, right=10pt, top=10pt, bottom=10pt,
    #1
}

% 2. Quote/Translation Box
\newtcolorbox{quotebox}[1][]{
    enhanced, breakable,
    frame hidden,
    colback=lightgray,
    borderline west={3pt}{0pt}{emerald},
    left=15pt, right=10pt, top=10pt, bottom=10pt,
    sharp corners,
    #1
}

% 3. AI-generated notice box
\newtcolorbox{warnbox}[1][]{
    enhanced, breakable,
    colback=red!4,
    colframe=darkred,
    boxrule=0.8pt,
    arc=3pt,
    left=10pt, right=10pt, top=10pt, bottom=10pt,
    #1
}

% Commands
\newcommand{\arab}[1]{\foreignlanguage{arabic}{#1}}
\newcommand{\kib}[1]{\textcolor{emerald}{\textbf{#1}}}

% Latin (English) fallback: Noto Serif Bengali has NO Latin glyphs
\newfontfamily\latfont[Renderer=HarfBuzz]{Noto Serif}
\newcommand{\lat}[1]{{\latfont #1}}

\setlength{\parskip}{5pt}
\setlength{\parindent}{0pt}

% Fix for missing bullet in Noto Serif Bengali
\renewcommand{\labelitemi}{$\bullet$}



\begin{document}
\begin{center}{\Large\bfseries\color{emerald} \lat{07}-কাউমা-দোয়া}\end{center}
\vspace{1em}
\section*{০৭ — কাউমা (রুকু থেকে দাঁড়ানো): \textit{\lat{Sami'Allahu} + \lat{Rabbana} \lat{Lakal} \lat{Hamd}}}

\begin{warnbox}
 \textbf{গুরুত্বপূর্ণ নোটিশ:} এই কন্টেন্টটি \lat{AI} দিয়ে প্রস্তুত। কেবল সহীহ/হাসান হাদিস রাখা হয়েছে। আলেম-যাচাইকৃত নয়।
\end{warnbox}

\medskip\hrule\medskip

\subsection*{১. সূত্র কার্ড}

\begin{table}[h]\centering\small
\begin{tabular}{ll}
বিষয় & বিবরণ \\
\hline
\textbf{বিষয়} & রুকু থেকে দাঁড়ানোর সময় ও দাঁড়ানোর পর বলা বাক্য \\
\textbf{সূত্র} & সহীহ বুখারি ৭৯৫, সহীহ বুখারি ৭৮৯ \\
\textbf{মর্যাদা} & \textbf{সুন্নাতে মুআক্কাদাহ} \\
\hline
\end{tabular}
\end{table}

\medskip\hrule\medskip

\subsection*{২. মূল আরবি}

\subsubsection*{\textbf{ইমাম বা একাই দাঁড়ানোর সময়:}}
\begin{center}\arab{سَمِيعَ اللَّهُ لِمَنْ حَمِدَهُ}\end{center}

\subsubsection*{\textbf{মুক্তাদি (জমাত) দাঁড়ানোর পর:}}
\begin{center}\arab{رَبَّنَا لَكَ الْحَمْدُ}\end{center}

\medskip\hrule\medskip

\subsection*{৩. বাংলা উচ্চারণ}

\textbf{সামিআল্লাহু লিমান হামিদাহ}
\textbf{রাব্বানা লাকাল হামদ}

\medskip\hrule\medskip

\subsection*{৪. শব্দ-শব্দ অর্থ}

\begin{table}[h]\centering\small
\begin{tabular}{lll}
আরবি & বাংলা & \lat{English} \\
\hline
\textbf{\arab{سَمِيعَ}} & সবকিছু শোনেন & \textit{\lat{Allah} \lat{hears}} \\
\textbf{\arab{اللَّهُ}} & আল্লাহ & \textit{\lat{Allah}} \\
\textbf{\arab{لِمَنْ}} & যে & \textit{\lat{who}} \\
\textbf{\arab{حَمِدَهُ}} & তাঁর প্রশংসা করে & \textit{\lat{praises} \lat{Him}} \\
\textbf{\arab{رَبَّنَا}} & হে আমাদের রব & \textit{\lat{Our} \lat{Lord}} \\
\textbf{\arab{لَكَ}} & তোমারই & \textit{\lat{for} \lat{You}} \\
\textbf{\arab{الْحَمْدُ}} & সকল প্রশংসা & \textit{\lat{all} \lat{praise}} \\
\hline
\end{tabular}
\end{table}

\medskip\hrule\medskip

\subsection*{৫. পূর্ণ বাংলা অনুবাদ}

\textbf{ইমাম:} \textbf{\lat{“}আল্লাহ তাঁর প্রশংসাকারীকে শোনেন।\lat{”}}

\textbf{মুক্তাদি:} \textbf{\lat{“}হে আমাদের রব! সকল প্রশংসা তোমারই।\lat{”}}

\medskip\hrule\medskip

\subsection*{৬. সংক্ষিপ্ত ব্যাখ্যা}

\subsubsection*{\textbf{"সামিআল্লাহু" — আল্লাহ শোনেন}}
রুকু থেকে দাঁড়ানোর সময় ইমাম বলেন: "সামিআল্লাহু লিমান হামিদাহ" — অর্থ: \textbf{"আল্লাহ তাঁর প্রশংসাকারীকে শোনেন।"} এটি \textbf{\lat{acknowledgment} (স্বীকৃতি)}: আমরা প্রশংসা করছি, আল্লাহ শুনছেন।

\subsubsection*{\textbf{"রাব্বানা লাকাল হামদ" — প্রশংসা তোমারই}}
মুক্তাদিরা এই বাক্যটি বলেন: \textbf{"হে আমাদের রব! সকল প্রশংসা তোমারই।"} এটি \textbf{\lat{hamd} (হামদ)}-এর পূর্ণতা — প্রশংসার অধিকারী শুধু আল্লাহ।

\subsubsection*{\textbf{কীভাবে দাঁড়াবেন?}}
\begin{itemize}
\item \textbf{ধীরে} দাঁড়াবেন — তাড়াহুড়ো নয়।
\item প্রতিটি \textbf{\lat{vertebra} (মেরুদণ্ড)} একে একে \textbf{\lat{stacked} (সাজানো)} হোক।
\item দাঁড়ানোর পর \textbf{সুস্থভাবে দাঁড়িয়ে} থাকুন — "ইতমিনান" (প্রত্যাশা)।
\end{itemize}
\medskip\hrule\medskip

\subsection*{৭. খুশুর জন্য মূল পয়েন্টস}

\begin{table}[h]\centering\small
\begin{tabular}{ll}
পয়েন্ট & কীভাবে প্রয়োগ করবেন \\
\hline
\textbf{ধীর গতি} & রুকু থেকে দাঁড়ানো — "আল্লাহু আকবার" বলে ধীরে উঠুন \\
\textbf{ইতমিনান} & দাঁড়ানোর পর কিছুক্ষণ স্থির হোন — "কামা ইলতাজাহাত" (ইতমিনান) \\
\textbf{প্রশংসা} & "রাব্বানা লাকাল হামদ" — হৃদয়ে অনুভব করুন \\
\textbf{নজর} & চোখ কিবলার দিকে — ভুমির দিকে \\
\hline
\end{tabular}
\end{table}

\medskip\hrule\medskip

\subsection*{৮. সংশ্লিষ্ট হাদিস}

\begin{table}[h]\centering\small
\begin{tabular}{ll}
হাদিস & গ্রন্থ \\
\hline
\textbf{\lat{“}যখন ইমাম 'সামিআল্লাহু লিমান হামিদাহ' বলেন, তোমরা 'রাব্বানা লাকাল হামদ' বলো।\lat{”}} & সহীহ বুখারি ৭৮৯ \\
\textbf{\lat{“}তোমাদের প্রত্যেকে যখন রুকু থেকে দাঁড়াবে, তখন দাঁড়িয়ে যাওয়ার পর ইতমিনান করুক।\lat{”}} & সহীহ বুখারি ৭৯৫ \\
\hline
\end{tabular}
\end{table}

\medskip\hrule\medskip

\textit{শেষ আপডেট: আগস্ট ২০২৬}


\end{document}
